%! Comparing Studies and Choosing a Method — Unit 1, Lesson 1.7 — Guided Notes & Practice (Answer Key)
% THE in-class document, in the MAIN IDEAS / NOTES shape (LESSON_SHAPE.md §1, ported from AP
% Statistics 2026-09-12): the vocabulary box, the hook, then ONE two-column guidednotes table.
% Each row is one idea — a label on the left; on the right one or two COMPLETE printed
% sentences, ONE large pre-drawn display the student reads or names, then a bold prompt and a
% two-across grid of numbered problems with 1.2–2.4cm of writing room. Problems run 1..19.
% DENSITY: a \blank{} only where a word IS the answer (the six cells of the side-by-side
% table in row 1 — the one \wordbank strip sits above it); no mid-sentence blanks; every
% display is pre-drawn; the page belongs to the student's pen.
%   I DO   : the teacher reads the definition, marks up the display, works the first problem
%            of each grid (2, 6, 10, 12).
%   WE DO  : the class works the rest of each grid, students holding the pen; the last row,
%            Guided Practice (16–19), is worked entirely together on Harbor Point instead of
%            Riverbend.
%   YOU DO : nothing on this page — ../homework/ IS the individual practice, started in the
%            last ten minutes of the period. No independent set, no reflection box, no
%            objectives box (the cover carries the targets).
% THE TRAP is problem 8 (the 45-minute bus ride: "unethical" is the wrong reason — it is
% IMPOSSIBLE, you cannot assign where a family lives).
% THE CRUX is problem 15 (row 4, "The Method"): the counselor SAT IN THE STUDY CENTER AND
% WATCHED who came — observation the METHOD — inside a study whose groups a generator built —
% an experiment, the DESIGN. The hook vote is re-taken there (PS.DC.3c, 3e). Problem 19 is
% the crux transferred to Harbor Point; problem 18 is the cannot-assign case transferred.
%
% THE DATA (verified in Python before authoring)
%   Riverbend, watched (1.5): 150 records, 60 users 39 passed -> 65%; 90 non-users 72 -> 80%; gap 15
%   Riverbend, assigned (1.6): 120 volunteers 60/60; 75% vs 55%; gap 20
%   Harbor Point (Guided Practice): 1.5 survey 200 members 70% vs 45% (gap 25, members chose
%     their swim time); 1.6 experiment 90 volunteers 45/45, 60% vs 40% (gap 20); NEW: 120
%     volunteers 60/60, six weeks, 42/60 = 70% vs 30/60 = 50%, gap 20
%
% PAGE LOCKSTEP with notes_key/: the key is GENERATED from this file by templates/lesson/mkkey.py
% (spec: templates/lesson/examples/lesson07_notes_key.json) — every \blank{} becomes an \ans{},
% every \vterm{} a \vtermans{}{}, every \par\writelines{n} n \ansline{} calls, and the answer
% argument of every \pcell, \writespace and \labelbox is filled. Answer spaces are fixed
% heights, so the two files cannot drift. KEEP EVERY \pcell ON ONE SOURCE LINE — mkkey matches
% line by line.
% TABLE MECHANICS: inside a cell, `\\` ENDS THE ROW — break lines with \par. A row cannot break
% across pages: the figure gets its own sub-row (`\\` then `& ...`) and EACH GRID ROW is its own
% sub-row — close the probgrid, `\\ &`, reopen as probgrid* (no top rule).
\documentclass[12pt]{article}
\usepackage{saar-article}
\usepackage{saar-key}   % pulls in -boxes; do NOT also load -boxes
\newcommand{\TallMath}[1]{$\displaystyle #1\rule[-1.4em]{0pt}{3.2em}$}
% Word bank strip — ONLY above a table the student fills (row 1). Defined per document;
% the key inherits it and prints the same strip, so heights match.
\newcommand{\wordbank}[1]{\par\vspace{2pt}\noindent\fcolorbox{goldacc}{goldbg}{\parbox{\dimexpr\linewidth-2\fboxsep-2\fboxrule\relax}{\small\textbf{Word bank:} #1}}\par\vspace{4pt}}
% Vocabulary rows: a FIXED-HEIGHT row carrying the term label and OPEN writing space —
% no underline beside the term and no rule line (user correction 2026-09-06: the black
% answer line crowded the row; students write beside and below the term). The key fills
% exactly the same height, so the box cannot drift blank vs. keyed. saar-article's
% \termblank adds an inline blank and a rule line, so the pair is defined here (the key is
% generated from this file and inherits both). 2.3cm per row gives students three
% handwritten lines of room; keep key definitions to two lines.
\newcommand{\termrowheight}{1.7cm}
\newlength{\termrowinset}\setlength{\termrowinset}{7pt}
\newcommand{\vterm}[1]{%
  \par\noindent\begin{minipage}[t][\termrowheight][t]{\linewidth}%
    \vspace*{\termrowinset}%
    \textbf{\textcolor{cerulean}{#1:}}%
  \end{minipage}\par}
\newcommand{\vtermans}[2]{%
  \par\noindent\begin{minipage}[t][\termrowheight][t]{\linewidth}%
    \vspace*{\termrowinset}%
    \textbf{\textcolor{cerulean}{#1:}}\quad\ans{#2}%
  \end{minipage}\par}

\begin{document}

\pageheader{Unit 1, Lesson 1.7}{Guided Notes \& Practice --- Answer Key}
% NAMESTRIP: no \namedateperiod here — mirrors the blank. NO teachernote here —
% teacher prose lives in the lesson plan.

\vspace{0.15in}

\begin{vocabbox}
\small
Fill in each term \textbf{as we name it} in the notes below.
\vtermans{Controlled experiment}{A study in which the researcher assigns the units to groups by chance and imposes a treatment on at least one group. It can support a cause-and-effect conclusion.}
\vtermans{Observational study}{A study in which the researcher measures or asks without assigning anyone to a group --- the groups already exist. It can report an association only.}
\vtermans{Collection method}{How the data arrive --- survey, interview, focus group, observation, or content analysis. Chosen by what the question needs, not by the design.}
\vtermans{Observation}{A collection method: recording what people actually do, as it happens. Not a design --- an experiment can collect its data by observation.}
\vtermans{Content analysis}{A collection method: studying records or documents that already exist --- borrowing histories, minutes, past scores. Best for change over time.}
\end{vocabbox}

\boxguard[12]
\begin{hookbox}
\small
This quarter the Riverbend counselor let a random number generator split $120$ volunteers into
two groups of $60$. To collect her data she handed out no questionnaire --- she sat in the Study
Center three afternoons a week and \emph{watched}. A classmate says: \emph{``She just watched, so
this quarter's study was an observational study, the same as last quarter's.''}

\vspace{3pt}
\textbf{Is the classmate right?} Circle one: \quad agree \qquad disagree
\quad --- and say why you think so. We vote again at problem~15.
\par\ansline{Answers vary --- most agree: ``she only watched, so it must be observational.''}
\ansline{(Settled at problem 15: disagree --- the generator built the groups; an experiment.)}
\end{hookbox}

\small
\begin{guidednotes}
% ── ROW 1: the two designs, side by side ─────────────────────────────────────
\mainidea[Side by side]{Two Designs} &
An \textbf{observational study} measures or asks without assigning anyone to a group --- the
groups already exist. A \textbf{controlled experiment} assigns the groups by chance and imposes
a treatment. \textbf{One question sorts them: who decided which group each individual is in?}
\\
& {\centering
\begin{tikzpicture}
  % left panel: LAST QUARTER — watched 150 records, the students sorted themselves
  \fill[frost, rounded corners=5pt] (0,0) rectangle (5.4,3.3);
  \draw[cerulean, thick, rounded corners=5pt] (0,0) rectangle (5.4,3.3);
  \node[font=\scriptsize\bfseries\color{cerulean}, align=center, text width=4.9cm] at (2.7,2.95) {LAST QUARTER --- watched $150$ records};
  \fill[goldbg] (1.45,1.55) ellipse (1.0 and 0.6);
  \draw[goldacc, very thick] (1.45,1.55) ellipse (1.0 and 0.6);
  \node[font=\scriptsize\bfseries] at (1.45,1.72) {$60$ used it};
  \node[font=\tiny] at (1.45,1.35) {$65\%$ passed};
  \fill[goldbg] (3.85,1.55) ellipse (1.15 and 0.7);
  \draw[goldacc, very thick] (3.85,1.55) ellipse (1.15 and 0.7);
  \node[font=\scriptsize\bfseries] at (3.85,1.72) {$90$ did not};
  \node[font=\tiny] at (3.85,1.35) {$80\%$ passed};
  \node[font=\tiny\itshape, align=center, text width=4.9cm] at (2.7,0.42) {the students sorted themselves --- unequal groups};
  % right panel: THIS QUARTER — assigned 120 volunteers, chance sorted them
  \fill[frost, rounded corners=5pt] (5.9,0) rectangle (11.3,3.3);
  \draw[cerulean, thick, rounded corners=5pt] (5.9,0) rectangle (11.3,3.3);
  \node[font=\scriptsize\bfseries\color{cerulean}, align=center, text width=4.9cm] at (8.6,2.95) {THIS QUARTER --- assigned $120$ volunteers};
  \node[draw=cerulean, fill=white, rounded corners=2pt, font=\tiny\bfseries, inner sep=2pt] (gen) at (8.6,2.35) {random number generator};
  \fill[goldbg] (7.35,1.3) ellipse (1.05 and 0.6);
  \draw[goldacc, very thick] (7.35,1.3) ellipse (1.05 and 0.6);
  \node[font=\scriptsize\bfseries] at (7.35,1.47) {$60$ treated};
  \node[font=\tiny] at (7.35,1.1) {$75\%$ passed};
  \fill[goldbg] (9.85,1.3) ellipse (1.05 and 0.6);
  \draw[goldacc, very thick] (9.85,1.3) ellipse (1.05 and 0.6);
  \node[font=\scriptsize\bfseries] at (9.85,1.47) {$60$ control};
  \node[font=\tiny] at (9.85,1.1) {$55\%$ passed};
  \draw[->, thick, cerulean] (gen.south) -- (7.35,1.92);
  \draw[->, thick, cerulean] (gen.south) -- (9.85,1.92);
  \node[font=\tiny\itshape, align=center, text width=4.9cm] at (8.6,0.42) {chance sorted them --- equal groups, nothing else changed};
\end{tikzpicture}\par\vspace{4pt}
Left: a(n) \labelbox{4.0cm}{observational study}\qquad Right: a(n) \labelbox{2.6cm}{experiment}\par}
\\
& \textbf{1.} Complete the side-by-side table. The first row is given --- it produces every other row.
\wordbank{yes \;\textbullet\; no \;\textbullet\; not necessarily \;\textbullet\; an association \;\textbullet\; cause and effect}
\footnotesize\renewcommand{\arraystretch}{1.05}
\begin{tabularx}{\linewidth}{@{} X p{3.3cm} p{3.3cm} @{}}
\toprule
\textbf{Ask this} & \textbf{Observational study} & \textbf{Controlled experiment} \\
\midrule
Who decides which group each individual is in? & the individuals themselves & chance --- a generator \\
Does the researcher impose a treatment? & \ans{no} & \ans{yes} \\
Are the groups alike before the study starts? & \ans{not necessarily} & \ans{yes} \\
Strongest conclusion it earns & \ans{an association} & \ans{cause and effect} \\
\bottomrule
\end{tabularx}\small
\\
& \notesprompt{Write the strongest sentence each Riverbend study earns --- the verb is the whole answer.}
\begin{probgrid}{|Y|Y|}
\pcell{2}{Watched: $65\%$ of the $60$ users passed, against $80\%$ of the $90$ non-users.}{1.5cm}{Using the Study Center is associated with a lower pass rate --- nothing stronger.} & \pcell{3}{Assigned: $75\%$ of the $60$ treated passed, against $55\%$ of the $60$ control.}{1.5cm}{The Study Center caused a higher pass rate --- chance built the groups.} \\ \hline
\end{probgrid}
\\
& \begin{probgrid}*{|Y|Y|}
\pcell{4}{Last quarter the Study Center looked \emph{harmful}, this quarter \emph{helpful}. Which result should she believe, and what earns it that trust?}{1.6cm}{This quarter's --- chance built two alike groups, so only the treatment differs.} & \pcell{5}{The warm-up's tell: one study's groups were $60$ and $90$, the other's $60$ and $60$. How do the sizes give away who built them?}{1.6cm}{Self-sorting leaves sizes 60/90; a generator splitting 120 makes 60/60.} \\ \hline
\end{probgrid}
\\ \hline
% ── ROW 2: when you cannot assign — three reasons, one test ──────────────────
\mainidea[When you cannot]{Experiment} &
To run an experiment you must be able to \textbf{assign} the treatment. Three things stop you:
\textbf{unethical} --- it would harm someone; \textbf{impossible} --- the treatment is something a
person already \emph{is} or \emph{has}; \textbf{impractical} --- too costly, too slow, or nobody
would agree. \textbf{The test: can you hand this out to a person?}
\\
& {\centering
\begin{tikzpicture}
  \node[draw=cerulean, fill=frost, thick, rounded corners=4pt, font=\small\bfseries, inner sep=5pt, align=center, text width=8.2cm] (q) at (5.3,2.75) {Can you hand the treatment out to a person --- by chance?};
  \node[draw=greenacc, fill=greenbg, thick, rounded corners=4pt, font=\scriptsize, inner sep=4pt, align=center, text width=4.2cm] (yes) at (2.6,1.05) {\textbf{YES --- run the experiment.}\\ Chance builds the groups, so they start out alike.\\ Strongest verb: \textbf{caused}.};
  \node[draw=redacc, fill=redbg, thick, rounded corners=4pt, font=\scriptsize, inner sep=4pt, align=center, text width=4.6cm] (no) at (8.1,1.05) {\textbf{NO --- unethical, impossible, or impractical.}\\ Run the best observational study you can.\\ Strongest verb: \textbf{is associated with}.};
  \draw[->, very thick, greenacc] (q.south) -- (yes.north) node[midway, left=2pt, font=\scriptsize\bfseries\color{greenacc}] {yes};
  \draw[->, very thick, redacc] (q.south) -- (no.north) node[midway, right=2pt, font=\scriptsize\bfseries\color{redacc}] {no};
  \node[font=\scriptsize\itshape\color{cerulean}, align=center, text width=10.6cm] at (5.3,-0.62) {``We could not run an experiment'' is the reason for a \textbf{weaker} verb --- never an excuse for a stronger one.};
\end{tikzpicture}\par}
\\
& \notesprompt{Can you assign it? If not, say which reason --- and what you run instead.}
\begin{probgrid}{|Y|Y|}
\pcell{6}{Does a $10$-minute silent reading block raise reading scores?}{1.6cm}{Yes --- you can hand a reading block out by chance. Run the experiment.} & \pcell{7}{Do students who smoke have worse lung function at age $30$?}{1.6cm}{No --- unethical: nobody may be assigned to smoke. Run an observational study.} \\ \hline
\end{probgrid}
\\
& \begin{probgrid}*{|Y|Y|}
\pcell{8}{Do students with bus rides over $45$ minutes score lower on the state test? A classmate says \emph{``unethical.''} Is that the reason?}{1.6cm}{No --- impossible, not unethical: you cannot assign where a family lives.} & \pcell{9}{The long riders \emph{do} score lower. May the report say the ride \emph{lowers} scores? Give the verb it may use, and why.}{1.6cm}{No. The groups built themselves --- the verb is \emph{is associated with}.} \\ \hline
\end{probgrid}
\\ \hline
% ── ROW 3: planning an experiment through the cycle ──────────────────────────
\mainidea[Planning an experiment]{The Cycle} &
The four stages of the \textbf{statistical cycle} plan an experiment too. Only stage two
changes: \stepnum{2}\ \textbf{Collect} now asks who the \emph{units} are and how chance
\textbf{assigns} them to groups --- comparison, randomization, replication, and control are all
checked here. The counselor's next question: \emph{does a weekly one-on-one check-in raise the
share of students who turn in every assignment?}
\\
& {\centering
\begin{tikzpicture}
  \node[draw=cerulean, fill=frost, thick, rounded corners=4pt, minimum width=2.3cm, minimum height=0.9cm, font=\small\bfseries\color{cerulean}] (f) at (1.2,0) {Formulate};
  \node[draw=goldacc, fill=goldbg, very thick, rounded corners=4pt, minimum width=2.3cm, minimum height=0.9cm, font=\small\bfseries\color{goldacc}] (c) at (3.9,0) {Collect};
  \node[draw=cerulean, fill=frost, thick, rounded corners=4pt, minimum width=2.3cm, minimum height=0.9cm, font=\small\bfseries\color{cerulean}] (r) at (6.6,0) {Represent};
  \node[draw=cerulean, fill=frost, thick, rounded corners=4pt, minimum width=2.7cm, minimum height=0.9cm, font=\small\bfseries\color{cerulean}] (a) at (9.5,0) {Analyze \& report};
  \draw[->, thick] (f) -- (c);
  \draw[->, thick] (c) -- (r);
  \draw[->, thick] (r) -- (a);
  \draw[->, thick, cerulean] (9.5,-0.45) -- (9.5,-0.85) -- (1.2,-0.85) -- (1.2,-0.45);
  \node[font=\tiny\itshape\color{cerulean}] at (5.35,-1.03) {the cycle runs again with the next question};
  \node[font=\scriptsize, align=center, text width=5.4cm] at (3.9,1.02) {Lesson 1.5: \emph{sample} the units\\ Lesson 1.7: \emph{assign} them by chance};
  \draw[->, goldacc, thick] (3.9,0.62) -- (3.9,0.47);
\end{tikzpicture}\par\vspace{5pt}
The stage where an experiment's plan differs from 1.5's: \labelbox{2.0cm}{Collect}\par}
\\
& \notesprompt{Plan the check-in experiment. Write the stage, not just its name.}
\begin{probgrid}{|Y|Y|}
\pcell{10}{\textbf{Collect.} Name the units, say what the generator does, and name the two groups.}{2.0cm}{Units: student volunteers. A generator sends half to a weekly check-in, half to none.} & \pcell{11}{\textbf{Analyze \& report.} What two numbers do you compare, and which verb may the report use? Why that verb?}{2.0cm}{Compare each group's percent turning everything in. Verb: \emph{caused} --- chance built the groups.} \\ \hline
\end{probgrid}
\\ \hline
% ── ROW 4: THE CRUX — five methods, and the method is not the design ─────────
\mainidea[Five ways to collect]{The Method} &
Once the design is set, you still choose \emph{how the data arrive}. A \textbf{collection
method} is how each answer gets recorded: \textbf{observation} records what people actually
\emph{do}, \textbf{content analysis} studies records that already exist. An \textbf{observational
study} is a \emph{design} --- nobody assigned the groups.
\\
& {\centering
\begin{tikzpicture}
  \foreach \i/\cname/\cdesc in {
      0/Survey/{the same fixed\\questions to\\many people},
      1/Interview/{depth from one\\person --- the\\reasons, the \emph{why}},
      2/Focus group/{a small group\\building on each\\other's answers},
      3/Observation/{record what\\people actually\\\emph{do}},
      4/{Content\\analysis}/{records that\\already exist ---\\change over time}} {
    \fill[frost, rounded corners=4pt] ({\i*2.25},0) rectangle ({\i*2.25+2.1},1.85);
    \draw[cerulean, thick, rounded corners=4pt] ({\i*2.25},0) rectangle ({\i*2.25+2.1},1.85);
    \node[font=\footnotesize\bfseries\color{cerulean}, align=center, text width=1.95cm] at ({\i*2.25+1.05},1.48) {\cname};
    \node[font=\tiny, align=center, text width=1.95cm] at ({\i*2.25+1.05},0.62) {\cdesc};
  }
\end{tikzpicture}\par}
\\
& \fcolorbox{cerulean}{frost}{\parbox{\dimexpr\linewidth-2\fboxsep-2\fboxrule\relax}{\centering
\textbf{Two questions, two answers.} \emph{Who built the groups?} decides the design. \emph{How were the answers recorded?} decides the method. An experiment can collect its data by observation.}}

\vspace{3pt}
\textbf{In your own words:} what decides the design, and what decides the method?
\writespace{1.2cm}{Who built the groups decides the design; how the answers were recorded decides the method.}
\\
& \notesprompt{Choose the method the question needs --- then settle the hook vote.}
\begin{probgrid}{|Y|Y|}
\pcell{12}{The library wants five years of borrowing records summarized.}{1.2cm}{Content analysis --- the records already exist.} & \pcell{13}{Why did four students stop coming to the Study Center? Name the method, and why a written survey is a poor fit.}{1.5cm}{Interviews --- \emph{why} needs follow-up; a fixed survey cannot.} \\ \hline
\end{probgrid}
\\
& \begin{probgrid}*{|Y|Y|}
\pcell{14}{Does the cafeteria's fruit actually get \emph{taken} --- not just reported taken?}{1.2cm}{Observation --- record what students actually do.} & \pcell{15}{\textbf{Vote again:} agree \; disagree. She sat and watched who came. Name the \emph{method} and the \emph{design}. Does watching make it an observational study?}{1.9cm}{Disagree. Method: observation; design: experiment --- a generator built the groups.} \\ \hline
\end{probgrid}
\\ \hline
% ── ROW 5: GUIDED PRACTICE — we do, together, a fresh context ────────────────
\mainidea[Guided practice]{Harbor Point} &
The Harbor Point pool director has already run the two studies below on her $1{,}500$ members.
\textbf{Now the board asks:} does a weekly water-aerobics class reduce reported joint pain? She
takes $120$ volunteers; a generator sends $60$ to the class and $60$ to no class for six weeks.
At the end, $42$ of the class group and $30$ of the no-class group report less joint pain.
\\
& {\centering
\begin{tikzpicture}
  \foreach \i/\ttl/\body in {
    0/{Lesson 1.5 --- surveyed $200$}/{morning swimmers $70\%$\\evening swimmers $45\%$\\[3pt]\emph{members chose their time}\\gap: $25$ points},
    1/{Lesson 1.6 --- assigned $90$}/{generator: $45$ / $45$\\$60\%$ vs.\ $40\%$\\[3pt]\emph{chance built the groups}\\gap: $20$ points},
    2/{NEW --- assigned $120$}/{generator: $60$ / $60$\\class vs.\ no class, six weeks\\[3pt]$42$ vs.\ $30$ report\\less joint pain}} {
    \fill[frost, rounded corners=5pt] ({\i*3.75},0) rectangle ({\i*3.75+3.5},2.95);
    \draw[cerulean, thick, rounded corners=5pt] ({\i*3.75},0) rectangle ({\i*3.75+3.5},2.95);
    \node[font=\scriptsize\bfseries\color{cerulean}, align=center, text width=3.3cm] at ({\i*3.75+1.75},2.62) {\ttl};
    \node[font=\scriptsize, align=center, text width=3.3cm] at ({\i*3.75+1.75},1.25) {\body};
  }
\end{tikzpicture}\par}
\\
& \notesprompt{We work these together.}
\begin{probgrid}{|Y|Y|}
\pcell{16}{Label the two earlier studies by design, and say who built the groups in each.}{1.7cm}{1.5: observational (members chose their time). 1.6: experiment (generator built 45/45).} & \pcell{17}{Compute the two rates in the water-aerobics study and the gap, then write the strongest sentence it earns.}{1.9cm}{$42/60 = 70\%$ vs.\ $30/60 = 50\%$, gap $20$ points. The class caused less joint pain.} \\ \hline
\end{probgrid}
\\
& \begin{probgrid}*{|Y|Y|}
\pcell{18}{Next question: \emph{do members who already have arthritis sleep worse?} Can she assign it? Which reason? What does she run, with which verb?}{2.0cm}{No --- impossible: nobody can hand out arthritis. Observational study; verb \emph{is associated with}.} & \pcell{19}{She measured the pain by asking all $120$ the same three questions; she also wants to know \emph{why} some quit. Name both methods --- does either change the design?}{2.0cm}{Pain: a survey. Why they quit: interviews. Neither changes the design.} \\ \hline
\end{probgrid}
\\ \hline
\end{guidednotes}

% ── THE NOTES END AT GUIDED PRACTICE ─────────────────────────────────────────
% There is deliberately no "you do" here: ../homework/ is the individual practice,
% started in class in the last ten minutes while the teacher circulates.

\end{document}
